\begin{abstractind}

\textit{Penelitian ini bertujuan untuk menerapkan algoritma machine learning, yaitu K-Nearest Neighbors (KNN), dalam memprediksi kualitas pengguna pada aplikasi e-learning berbasis web. Dengan menggunakan dataset dummy yang terdiri dari fitur seperti waktu akses rata-rata, jumlah login per minggu, durasi video yang ditonton, dan rating pengguna, model KNN dilatih untuk mengklasifikasikan pengguna ke dalam kategori "Baik", "Cukup", dan "Buruk". Hasil evaluasi menunjukkan bahwa akurasi model mencapai 87\% dengan nilai k=5. Model ini dapat diintegrasikan ke dalam sistem rekomendasi untuk meningkatkan kualitas layanan e-learning.}

\bigskip
\noindent
\textbf{Kata kunci :} KNN, prediksi kualitas pengguna, e-learning, \textit{machine learning}, dan \textit{web-based application}
\end{abstractind}